\section{Problem formulation and research plan}
\label{sec:problem}

We consider the passive scalar advection--diffusion equation
\begin{equation}
\partial_t \theta + u^\omega \cdot \nabla \theta = D \Delta \theta,
\qquad \nabla \cdot u^\omega = 0,
\qquad \theta|_{t=0} = \theta_0,
\label{eq:scalar}
\end{equation}
on the periodic domain $\mathbb{T}^d$, with $d\in\{2,3\}$, scalar diffusivity $D>0$, and an autonomous random divergence-free velocity field $u^\omega=u^\omega(x)$ indexed by a random seed $\omega \in \Omega$. Throughout, we assume
\[
\langle \theta_0 \rangle := \frac{1}{|\mathbb{T}^d|}\int_{\mathbb{T}^d}\theta_0(x)\,dx = 0,
\]
so that the long-time limit of $\theta$ is trivial and all dissipation and regularity measurements must be made over finite, dynamically relevant time intervals.

The central goal is to determine how the roughness of the velocity field controls the onset of anomalous scalar dissipation in the limit $D\to0$, and how this transition is related to the spatial regularity of the scalar field itself. We are particularly interested in identifying a critical velocity roughness threshold $\alpha_c$ separating a regime with vanishing dissipation from a regime with persistent dissipation, and in testing whether the Yaglom-type relation
\begin{equation}
\alpha + 2\beta = 1
\label{eq:yaglom}
\end{equation}
correctly predicts the scalar H\"older exponent $\beta$.

\subsection{Velocity classes and notation}

We study three classes of autonomous incompressible velocity fields, all parameterized by a prescribed roughness exponent $\alpha\in[-1,1]$. Here, $C^\alpha$ is understood in the H\"older--Besov sense when $\alpha<0$.

\paragraph{(i) Two-dimensional stream-function class.}
When $d=2$, the velocity is generated by a random stream function $\psi^\omega$ via
\begin{equation}
u^\omega = \nabla^\perp \psi^\omega
:= \bigl(-\partial_2 \psi^\omega,\partial_1 \psi^\omega\bigr).
\label{eq:stream}
\end{equation}

\paragraph{(ii) Three-dimensional Leray-projected class.}
When $d=3$, we consider a general random vector field $v^\omega$ and define
\begin{equation}
u^\omega = \mathbb{P}v^\omega,
\qquad
\mathbb{P} := I - \nabla \Delta^{-1}\nabla\cdot,
\label{eq:leray}
\end{equation}
where $\mathbb{P}$ is the Helmholtz--Leray projection onto divergence-free vector fields.

\paragraph{(iii) Three-dimensional Clebsch class.}
Again in $d=3$, we consider structured incompressible velocities of Clebsch form
\begin{equation}
u^\omega = \nabla \phi_1^\omega \times \nabla \phi_2^\omega
\label{eq:clebsch}
\end{equation}
for random scalar potentials $\phi_1^\omega$ and $\phi_2^\omega$.

For notational convenience, we denote these three ensembles by
\[
\mathcal{C}\in
\left\{
\mathrm{2D\text{-}SF},
\mathrm{3D\text{-}Leray},
\mathrm{3D\text{-}Clebsch}
\right\}.
\]

For each class $\mathcal{C}$, exponent $\alpha$, diffusivity $D$, and seed $\omega$, let $\theta^{\omega}_{\mathcal{C},\alpha,D}$ denote the solution of \eqref{eq:scalar}. We then define the finite-time integrated scalar dissipation
\begin{equation}
\chi^{\omega}_{\mathcal{C},\alpha}(D,T)
:=
D \int_0^T \int_{\mathbb{T}^d}
|\nabla \theta^{\omega}_{\mathcal{C},\alpha,D}(x,t)|^2
\,dx\,dt,
\label{eq:chi}
\end{equation}
and its ensemble average
\begin{equation}
\overline{\chi}_{\mathcal{C},\alpha}(D,T)
:=
\mathbb{E}_\omega
\bigl[
\chi^{\omega}_{\mathcal{C},\alpha}(D,T)
\bigr].
\label{eq:chibar}
\end{equation}

We say that \emph{anomalous scalar dissipation} is observed at time horizon $T$ if
\begin{equation}
\liminf_{D\to0}\,
\overline{\chi}_{\mathcal{C},\alpha}(D,T) > 0.
\label{eq:anomalous}
\end{equation}
If instead $\overline{\chi}_{\mathcal{C},\alpha}(D,T)\to0$ as $D\to0$, then there is no anomalous dissipation on that observation window.

To quantify scalar regularity, we let $\beta=\beta_{\mathcal{C}}(\alpha)$ denote the measured spatial H\"older exponent of $\theta$, estimated from either a modulus-of-continuity fit or, equivalently, from second-order structure functions
\begin{equation}
S_2^\theta(r;t)
:=
\mathbb{E}_\omega
\left[
\left\langle
|\theta(x+r,t)-\theta(x,t)|^2
\right\rangle
\right]
\sim |r|^{2\beta}
\label{eq:structure}
\end{equation}
over an appropriate range of scales. Here $\langle \cdot \rangle$ denotes spatial averaging over $\mathbb{T}^d$.

\subsection{Assumptions}

The investigation is carried out under the following assumptions.

\begin{enumerate}
    \item The velocity field is autonomous in time and divergence-free by construction in each ensemble.

    \item The roughness exponent $\alpha$ is prescribed a priori and varied over a broad range, from smooth fields ($\alpha=1$) to very rough fields ($\alpha\approx -1$).

    \item The scalar initial data $\theta_0$ is fixed across experiments, has zero spatial mean, and is sufficiently regular to ensure well-posedness of \eqref{eq:scalar} for every realization under consideration.

    \item For each velocity class, the large-scale amplitude and length scale are normalized uniformly in $\alpha$ and $\omega$ so that comparisons across ensembles are meaningful.

    \item Since $\theta(\cdot,t)$ relaxes to its mean value as $t\to\infty$, all diagnostics are performed on finite observation windows
    \begin{equation}
    T = c\,\tau_{\mathrm{eddy}},
    \qquad c = O(1),
    \label{eq:timewindow}
    \end{equation}
    where $\tau_{\mathrm{eddy}}$ is a characteristic large-eddy turnover time associated with the given velocity realization.

    \item Whenever the velocity construction involves a velocity-side viscosity or smoothing parameter $\nu$, we restrict attention to the regime
    \begin{equation}
    Sc = \frac{\nu}{D} < 1.
    \label{eq:schmidt}
    \end{equation}

    \item Statistical conclusions are drawn from ensemble averages over many independent random seeds $\omega$, together with sensitivity checks in the observation time $T$.
\end{enumerate}

\subsection{Objective}

The primary objective of this project is to determine, for each velocity class $\mathcal{C}$, the critical roughness threshold at which anomalous scalar dissipation first appears as $D\to0$. A second objective is to measure the scalar regularity exponent $\beta$ and test whether the Yaglom-type balance \eqref{eq:yaglom} correctly describes the relation between velocity roughness and scalar roughness in each regime. A third objective is to compare the transition mechanisms across the three incompressible velocity constructions and to determine whether the general three-dimensional Leray-projected class exhibits behavior that differs from the more structured stream-function and Clebsch classes.

More concretely, we seek a quantitative description of
\begin{equation}
\overline{\chi}_{\mathcal{C},\alpha}(D,T)
\sim D^{\gamma_{\mathcal{C}}(\alpha,T)}
\qquad \text{as } D\to0,
\label{eq:gamma}
\end{equation}
where the exponent $\gamma_{\mathcal{C}}(\alpha,T)$ serves as a diagnostic of dissipation behavior. In particular, $\gamma_{\mathcal{C}}(\alpha,T)=0$ is consistent with anomalous dissipation, while $\gamma_{\mathcal{C}}(\alpha,T)>0$ indicates decay of the integrated dissipation in the vanishing-diffusivity limit.

\subsection{Research questions}

The study is organized around the following questions.

\begin{enumerate}
    \item For each class $\mathcal{C}$, does there exist a critical exponent $\alpha_c$ such that anomalous dissipation is absent for $\alpha>\alpha_c$ and present for $\alpha<\alpha_c$?

    \item How does the scaling exponent $\gamma_{\mathcal{C}}(\alpha,T)$ in \eqref{eq:gamma} depend on the velocity roughness $\alpha$ and on the observation time $T$?

    \item Does the scalar H\"older exponent $\beta$ satisfy the Yaglom-type relation $\alpha+2\beta=1$, and if not, in which regimes does this relation fail or only hold approximately?

    \item Does the onset of anomalous dissipation coincide with the onset of Yaglom-type scaling, or do these transitions occur at different values of $\alpha$?

    \item To what extent does the answer depend on the geometric structure of the velocity field, especially when comparing the general three-dimensional Leray-projected ensemble with the more constrained Clebsch ensemble?
\end{enumerate}

\subsection{Methodology}

Our methodology is designed to isolate the effects of velocity roughness, diffusivity, and ensemble structure in a systematic way.

\begin{enumerate}
    \item \textbf{Parameter sweep in roughness.}
    For each velocity class $\mathcal{C}$, choose a broad grid of roughness exponents $\alpha\in[-1,1]$ and generate ensembles of autonomous random divergence-free velocity fields with the prescribed regularity.

    \item \textbf{Parameter sweep in diffusivity.}
    For each pair $(\mathcal{C},\alpha)$, solve \eqref{eq:scalar} for a range of scalar diffusivities $D\ll1$, while keeping the initial data and large-scale normalization fixed.

    \item \textbf{Ensemble averaging.}
    Repeat the computation for many independent seeds $\omega$ in order to estimate the sample average \eqref{eq:chibar} and quantify variability across realizations.

    \item \textbf{Finite-time dissipation diagnostics.}
    Compute the integrated dissipation $\chi^{\omega}_{\mathcal{C},\alpha}(D,T)$ in \eqref{eq:chi} over observation windows $T=c\,\tau_{\mathrm{eddy}}$ with several values of $c$, and fit the resulting dependence on $D$ to extract the effective exponent $\gamma_{\mathcal{C}}(\alpha,T)$.

    \item \textbf{Scalar regularity diagnostics.}
    Estimate the scalar exponent $\beta$ from structure functions of the form \eqref{eq:structure}, or from an equivalent modulus-of-continuity or wavelet-based estimator, applied consistently across all three ensembles.

    \item \textbf{Test of Yaglom-type scaling.}
    Compare the measured values of $\alpha$ and $\beta$ against the prediction \eqref{eq:yaglom}, and determine whether deviations from this relation correlate with the absence of anomalous dissipation.

    \item \textbf{Identification of the transition regime.}
    Assemble the measurements of $\overline{\chi}_{\mathcal{C},\alpha}(D,T)$ and $\beta$ into a phase diagram in the $(\alpha,D)$-plane, with the aim of locating the roughness threshold at which the dissipation behavior changes and comparing this threshold across the three velocity classes.
\end{enumerate}

\subsection{Summary of the problem}

In summary, the problem is to determine how the vanishing-diffusivity behavior of passive scalar dissipation depends on the roughness and structure of an autonomous incompressible velocity field. The key unknown is the transition from a non-anomalous regime to an anomalous one as the velocity roughness exponent $\alpha$ decreases, together with the corresponding behavior of the scalar H\"older exponent $\beta$. The comparison between the two-dimensional stream-function ensemble, the three-dimensional Clebsch ensemble, and the more general three-dimensional Leray-projected ensemble is intended to clarify which aspects of the transition are universal and which depend on the geometric mechanism used to enforce incompressibility.